\chapter{Kidney Function Tests}
\index{Kidney Function Tests}

\section*{Definition and Overview}
Kidney function tests are blood and urine tests that measure how well the kidneys are filtering waste and maintaining the body's balance. The key blood tests are creatinine, a waste product of muscle activity, and the estimated glomerular filtration rate (eGFR), which is calculated from creatinine, age, and sex and expresses kidney function as a percentage of normal. Urine tests look for protein and blood, which signal kidney damage. These tests are used to diagnose kidney disease, monitor its progression, adjust medicines, and screen people at risk, such as those with diabetes and high blood pressure. They are simple, inexpensive, and among the most informative tests in medicine.

\section*{Epidemiology}
Chronic kidney disease affects around one in ten people, yet most people with early disease are unaware of it, because it causes no symptoms until advanced. Kidney function tests are among the most commonly ordered blood tests, used in routine health checks, chronic disease management, and hospital care. Around a quarter of people with diabetes and a similar proportion with high blood pressure have some degree of kidney impairment. Early detection through testing allows treatment that slows progression and prevents the need for dialysis.

\section*{Aetiology and Pathophysiology}
The kidneys filter waste and excess fluid from the blood through millions of tiny filters called nephrons. Creatinine, produced at a steady rate from muscle, is removed by the kidneys, so its blood level rises as filtration falls. The eGFR estimates the filtration rate: an eGFR above 90 is normal, 60 to 89 mildly reduced, and below 60 for three months indicates chronic kidney disease, with below 15 indicating kidney failure. Protein in the urine, especially albumin, signals damage to the filters and is a strong predictor of progression and heart risk. Blood, from inflammation, stones, or other damage, is also detected. The tests track damage from diabetes, high blood pressure, and other causes.

\section*{Clinical Features}
\begin{itemize}
    \item Blood tests: creatinine and eGFR, drawn from the arm
    \item Urine tests: a dipstick for protein and blood, or a laboratory test of albumin
    \item A single abnormal result is repeated, as it can be affected by illness, exercise, and medicines
    \item eGFR is reported as a number, with normal above 90
    \item Kidney disease stages run from 1 (normal function with signs of damage) to 5 (failure)
    \item Protein in the urine is reported with a score such as ACR
    \item Results guide the need for further tests, specialist referral, and medicine adjustment
\end{itemize}

\section*{Differential Diagnosis}
An abnormal kidney test result must be interpreted in context.
\begin{itemize}
    \item \textbf{Acute changes:} dehydration, infection, or acute kidney injury
    \item \textbf{Chronic kidney disease:} consistently reduced function over months
    \item \textbf{Medicines:} some drugs raise creatinine without damaging the kidney
    \item \textbf{Muscle mass:} creatinine is lower with low muscle mass and higher in large people
    \item \textbf{Protein in the urine:} transient with exercise or infection, or persistent with disease
    \item \textbf{Other conditions:} heart failure and liver disease affect the results
\end{itemize}

\section*{When to Seek Medical Care}
Kidney function testing is recommended for people with diabetes, high blood pressure, heart disease, a family history of kidney disease, or symptoms such as swelling, foamy urine, or blood in the urine, and it is part of routine health checks. An abnormal result should be discussed with a doctor, and urgent care is needed for a sudden drop in urine, severe swelling, breathlessness, or confusion.

\textbf{Contact your local emergency services for:}
\begin{itemize}
    \item Difficulty breathing or severe swelling
    \item Passing very little or no urine
    \item Confusion, drowsiness, or a seizure
    \item Vomiting with an inability to keep fluids down
    \item Chest pain or collapse
\end{itemize}

\section*{Diagnosis and Investigations}
\begin{itemize}
    \item \textbf{Blood creatinine and eGFR:} the core tests
    \item \textbf{Urine albumin:} the albumin-to-creatinine ratio (ACR)
    \item \textbf{Urine dipstick:} for protein, blood, and glucose
    \item \textbf{Electrolytes:} potassium, sodium, and other markers
    \item \textbf{Ultrasound:} assessment of kidney size and structure
    \item \textbf{Repeat testing:} to confirm persistent changes
    \item \textbf{Referral:} to a kidney specialist for significant or progressive disease
\end{itemize}

\section*{Management}
\begin{itemize}
    \item \textbf{Controlling the cause:} blood pressure, diabetes, and cholesterol
    \item \textbf{Medicines:} ACE inhibitors or ARBs, which protect the kidneys
    \item \textbf{Lifestyle:} healthy diet, low salt, weight management, and stopping smoking
    \item \textbf{Medicine review:} adjusting doses and avoiding medicines that harm the kidneys
    \item \textbf{Monitoring:} regular repeat testing to track trends
    \item \textbf{Treatment of complications:} anaemia, bone disease, and heart risk
    \item \textbf{Specialist care:} for advanced or progressive disease
    \item \textbf{Planning:} discussion of dialysis and transplant options when needed
\end{itemize}

\section*{Complications}
\begin{itemize}
    \item Progression of chronic kidney disease to kidney failure
    \item Heart disease and stroke, which are more common with kidney disease
    \item Anaemia, bone disease, and malnutrition
    \item Side effects of medicines
    \item Complications of over-treatment, including dehydration from fluid restriction
    \item Anxiety from abnormal results without proper explanation
    \item The need for dialysis and transplantation when disease progresses
\end{itemize}

\section*{Prognosis and Outcomes}
Kidney function tests enable early detection and treatment that slows or prevents progression. With good control of blood pressure and diabetes, many people with early kidney disease maintain stable function for years, and the decline to failure is prevented in most. The tests also guide the safe use of medicines and allow planning for the small proportion of people who do progress. Regular testing in at-risk people is one of the most effective preventive strategies in medicine.

\section*{Prevention and Self-Care}
\begin{itemize}
    \item Have kidney function tests if you have diabetes, high blood pressure, or risk factors
    \item Attend repeat testing as advised
    \item Control blood pressure and blood sugar
    \item Maintain a healthy weight and limit salt
    \item Do not smoke
    \item Drink enough fluid and avoid dehydration
    \item Use over-the-counter pain medicines with care
    \item Ask your doctor about your eGFR and urine results
    \item Seek urgent care for sudden symptoms
\end{itemize}

\section*{Sources and Further Reading}
The following reputable sources provide further information for healthcare students and patients seeking reliable, current advice about kidney function tests.
\begin{itemize}
    \item Pathology Tests Explained. \textit{Kidney function tests.} https://pathologytestsexplained.org.au/
    \item Kidney Health Australia. \textit{Kidney health checks.} https://kidney.org.au/
    \item Healthdirect Australia. \textit{Kidney function tests.} https://www.healthdirect.gov.au/
    \item The Royal Australian College of General Practitioners. \textit{Chronic kidney disease guidelines.} https://www.racgp.org.au/
    \item NHS. \textit{Chronic kidney disease tests.} https://www.nhs.uk/
    \item National Kidney Foundation. \textit{eGFR and kidney tests.} https://www.kidney.org/
\end{itemize}
